\chapter{Preparing the Recoil Run}

The measurement begins before any atom is kicked. It begins with the choice of carrier, the fixing of gauges, and
the writing of a tape clean enough that a later reader can replay the run without guessing what was meant. The
apparatus is merciless about this. A number without its provenance is not a measurement; it is a rumor with
digits. To measure \(\alpha\), the user must make the recoil countable, make the gauge explicit, and make every
choice that can move the answer visible before the answer is known.

\section{Choose the Carrier}

Choose an atom \(X\). The carrier must be heavy enough to handle, cold enough to keep its velocity narrow, simple
enough that its mass is known, and optically cooperative enough that light can kick it in a repeatable way.
Rubidium and cesium are the usual choices because their atomic masses, laser transitions, cooling methods, and
interferometer behavior have been worked hard by metrology. The apparatus does not care about the romance of the
species. It asks only whether the carrier makes the formula auditable:
\[
  \alpha =
  \left(
  \frac{2R_\infty}{c}\,
  \frac{A_{\rm r}(X)}{A_{\rm r}(e)}\,
  \frac{h}{m_X}
  \right)^{1/2}.
\]
If \(A_{\rm r}(X)\) is poorly known, the carrier contaminates the measurement no matter how beautiful the recoil
fringes look. If the atom cannot be cooled and addressed cleanly, the recoil tape fills with haze. A good carrier
is not the one that gives a convenient signal once; it is the one whose mass, transitions, systematics, and
repeatability can survive the ledger.

Prepare the carrier as a narrow packet. Cool the atoms, select a velocity class, and record the selection as a
tange operation: this cloud, this internal state, this launch, this geometry. The apparatus must know which atom
was asked. Then repeat the preparation until the selected packet is not an anecdote. Repetition turns the packet
from a happening into a measurement object. Only then is the carrier ready to be kicked.

\section{Make the Recoil Countable}

The kick is a photon momentum transferred to the atom. In a simple one-photon picture the recoil velocity obeys
\[
  v_r = \frac{h}{m_X\lambda}.
\]
In an atom interferometer the light usually acts through an effective wave vector \(k_{\rm eff}\), and the
recoil is read as a phase or frequency shift. A useful bookkeeping form is
\[
  \omega_r = \frac{\hbar k_{\rm eff}^2}{2m_X},
  \qquad
  \frac{h}{m_X} = \frac{4\pi\omega_r}{k_{\rm eff}^2}.
\]
The exact experimental pulse sequence may be Ramsey-Borde, Bloch oscillation, Bragg diffraction, or another
variant. The apparatus does not require one sacred layout. It requires that the layout convert photon momentum
into a counted phase shift whose scale is known.

That conversion is where the old machine reappears. The photon kicks are counted. The phase is read against a
clock. The wavelength and geometry are gauges. The fringe contrast carries noise. The atom's path through the
pulses writes a trace. The user does not merely take a final fringe and write down a number. The user records the
whole tape: pulse times, laser frequencies, effective wave vector, magnetic field setting, atom temperature,
cloud position, launch velocity, pulse power, alignment, and the sign convention for which recoil direction
counts as positive. Each item is either a tanged label or a funged count. Nothing important is allowed to float
as laboratory mood.

The raw observable should be blind if the laboratory can bear the discipline. Hide an offset in the phase or
frequency analysis so the operator cannot steer the apparatus toward the expected \(137\). The hidden offset is
not a trick; it is a fence. It keeps the reader's desire from becoming a force on the gauge. Only after the
systematic ledger is closed should the blind be opened and \(h/m_X\) revealed.

\section{Make the Gauge Honest}

The wavelength gauge must be traceable. The laser frequency fixes \(\lambda\) through \(c=\lambda\nu\), and
\(c\) is exact only because the SI has made it exact. The apparatus still has to measure or lock \(\nu\), account
for refractive paths, and state whether \(k_{\rm eff}\) is the simple optical wave vector or a multi-photon
effective one. A small angular misalignment changes the projection of the momentum; a small frequency error
changes the scale of the kick. These are not footnotes. They are ways the hand can move the answer.

The time gauge must be equally visible. The interferometer phase is accumulated over pulse separations and hold
times. If a clock is used, its calibration belongs in the ledger. If the experiment extracts a frequency by
scanning a detuning, the scan step, fit model, line shape, and residuals belong in the ledger. The apparatus does
not accept "the fit found the center" as a primitive. It asks how the center was found, what was allowed to vary,
what was held fixed, and how the residuals behaved when the data were split into independent halves.

The environmental gauges must be recorded before they are corrected. Gravity gradients, magnetic fields,
wavefront curvature, AC Stark shifts, blackbody fields, atom density, and background gas collisions can all
change the recoil phase or the inferred wave vector. Some will be negligible; none may be invisible. To ignore a
small correction because it is small is to confuse size with grade. The grade of a correction is earned by a
bound, not by a shrug.

\section{The Prepared Tape}

At the end of preparation the user should have three tapes. The first is the raw measurement tape: shots, phases,
fringe outcomes, timestamps, and atom settings. The second is the gauge tape: wavelength, timing, geometry,
calibrations, and sign conventions. The third is the correction tape: every known effect that can move the
recoil, with a value, an uncertainty, and the evidence for both.

Only then should the apparatus be allowed to compute \(h/m_X\). Preparation is not ancillary labor. It is the
measurement before the measurement, the part that prevents the final number from becoming a lucky digit. A clean
recoil run is a run whose tape can be handed to an unfriendly reader and still speak.

\section{Entering the Run in Lean}

After the laboratory has produced the raw recoil value, enter it as a \(\texttt{Datum}\). The value is the
measured \(h/m_X\). The \(\texttt{sigma}\) field is its standard statistical uncertainty before systematic
corrections are added. The grade is \(\texttt{measured}\), because the atom cloud supplied this number.

\begin{quote}\small
\begin{verbatim}
def rawRecoil : Datum :=
  { value := 4.591359258e-9
    sigma := 0.000000006e-9
    grade := .measured }
\end{verbatim}
\end{quote}

The digits above are an example slot, not a command to use those digits. Replace the value and uncertainty with
the closed laboratory values for the chosen atom. Keep more digits than the final report will print. Rounding
belongs at the end, after the bridge and uncertainty propagation have been run.

Next enter the imported constants in one record. The speed of light has zero uncertainty in SI because it is a
defining constant; the Rydberg constant and relative atomic masses carry the uncertainties supplied by the
constants adjustment. The grade is \(\texttt{given}\), because these quantities do not come from the recoil run:

\begin{quote}\small
\begin{verbatim}
def rubidium87Constants : Constants :=
  { rydberg :=
      { value := 10973731.568157
        sigma := 0.000012
        grade := .given }
    relativeMassX :=
      { value := 86.909180520
        sigma := 0.000000010
        grade := .given }
    relativeMassElectron :=
      { value := 0.000548579909065
        sigma := 0.000000000016
        grade := .given }
    lightSpeed :=
      { value := 299792458.0
        sigma := 0.0
        grade := .given } }
\end{verbatim}
\end{quote}

This record is also where the chosen constants adjustment is pinned. Do not mix editions casually. A Rydberg
constant from one adjustment and a mass ratio from another may still produce a number, but it is no longer a
single audited computation. If correlations are later needed, this is the seam where the covariance information
enters.

The working habit is therefore:
\[
\begin{array}{c|l}
1 & \hbox{record the raw recoil as measured data}\\
2 & \hbox{record imported constants as given data}\\
3 & \hbox{record each systematic effect as a named correction}\\
4 & \hbox{close the recoil only after the correction ledger is frozen}\\
5 & \hbox{close \(\alpha\) from the corrected recoil and constants}\\
6 & \hbox{compare after the report object has been computed}
\end{array}
\]
Those six moves are the computation manual in miniature. The Lean file merely makes them hard to blur.
